\section{Introduction}
\label{sec:introduction}

Chiplet-based systems partition functions across dies and reconnect them within
a package. This decomposition can improve reuse and allow each function to use
an appropriate process, but it moves system correctness into the die-to-die
interface and package boundary \cite{li2020chiplet,lau2021chiplets}. UCIe
defines an open layered interface intended to support such heterogeneous
integration \cite{sharma2022ucie}. Within its physical-layer control, the Link
Training State Machine (LTSM) coordinates reset, sideband initialization,
mainband initialization and training, active operation, retraining, error, and
low-power behavior \cite{ucie20}.

An academic RTL model of this control is easy to overstate. Listing standard
states does not implement their physical procedures; a passing nominal
simulation does not test message ordering under backpressure; and an FPGA
synthesis report does not establish a working UCIe link. Reproducible negative
evidence is equally important: failed fits, malformed responses, pattern
failure, abort timing, and retained error behavior define the boundary of what
has actually been shown.

This paper presents a source-available, specification-traceable RTL and
verification framework for a UCIe 2.0-derived digital-control subset. The
synthesizable design combines a nine-state hierarchical controller, a bounded
retrying sideband sequencer, a 16-lane 23-bit LFSR measurement engine, retained
TRAINERROR classification, and a compact FPGA CSR wrapper. Its detailed case
study is the production DATATRAINCENTER1 path, in which START completion admits
pattern measurement, pattern failure retries without END, and a matching END
response gates advancement.

The paper makes three bounded contributions:

\begin{enumerate}
\item a modular controller architecture that integrates sideband transaction,
accepted-sample LFSR measurement, error retention, and compact observability
without presenting the result as a complete UCIe PHY;
\item an interface-derived UVM predictor and SVA layer that independently check
message order, per-lane LFSR state, strict threshold semantics, abort cause, and
state advancement; and
\item a reproducible evidence chain comprising directed suites, preserved
multi-seed campaigns, 180 production integrated trials across nine scenario
classes, and qualified Quartus results that retain both a failed wide-top fit
and a successful 119-pin CSR-wrapper fit.
\end{enumerate}

Section~\ref{sec:architecture} defines the implementation boundary and RTL.
Section~\ref{sec:verification} describes prediction, assertions, and scenario
accounting. Section~\ref{sec:evaluation} reports simulation and FPGA evidence.
Sections~\ref{sec:related} and \ref{sec:discussion} position the contribution
and state its limitations.
